Si tenemos un espacio vectorial de dimensión finita, sabemos que su espacio dual
es isomorfo a este, por lo que el doble dual también lo es. Sin embargo,
mientras el isomorfismo con el dual es algo más
\foreignlanguage{english}{ad-hoc}, el del doble dual se ve como algo más
fundamental, pues se obtiene <<de forma natural>> un isomorfismo que podríamos
llamar <<canónico>>, dado por $v\mapsto(f\mapsto f(v))$.

El hecho de considerar una cierta transformación, o una cierta operación, como
natural es algo común en distintas áreas de las matemáticas, y si bien el uso
del término suele ser informal, en la primera mitad del siglo XX se hicieron
esfuerzos por formalizarlo que, de hecho, fueron los que llevaron a la creación
de la teoría de categorías. En este capítulo exploramos este concepto,
basándonos principalmente en \cite[cap. 6]{joyofcats} y \cite[I.4 y
II.4--5]{maclane}.

\begin{definition}
  Dados dos morfismos $S,T:\cC\to\cD$, una \conc{transformación natural}
  $\tau:S\to T$, también escrita como
  \[\cC\natg{\downarrow\tau}{S}{T}\cD,\]
  es una función $\tau:\Ob{\cC}\to\Mor{\cD}$ que a cada objeto $c$ en $\cC$ le
  asigna un morfismo $\tau_c:Sc\to Tc$ de forma que, para todo morfismo
  $f:a\to b$ en $\cC$, la figura \ref{fig:natural} conmuta, esto es,
  $Tf\circ\tau_a=\tau_b\circ Sf$.
  \begin{figure}
    \centering
    \begin{diagram}
      \path (0,2) node(A){$a$} (2,2) node(SA){$Sa$} (4,2) node(TA){$Ta$};
      \path (0,0) node(B){$b$} (2,0) node(SB){$Sb$} (4,0) node(TB){$Tb$};
      \draw[->] (A) -- node[right]{$f$} (B);
      \draw[->] (SA) -- node[left]{$Sf$} (SB);
      \draw[->] (TA) -- node[right]{$Tf$} (TB);
      \draw[->] (SA) -- node[above]{$\tau_a$} (TA);
      \draw[->] (SB) -- node[below]{$\tau_b$} (TB);
    \end{diagram}
    \caption{Transformación natural.}
    \label{fig:natural}
  \end{figure}
\end{definition}

Si pensamos en los funtores $S$ y $T$ como imágenes en $\cD$ de los objetos y
los morfismos de $\cC$, una transformación natural nos da un conjunto de
morfismos de la imagen de $S$ a la de $T$ de forma que todos los cuadrados como
los de la figura \ref{fig:natural} conmutan.

También podemos pensar en la transformación natural como un morfismo
<<genérico>> en $\cD$, parametrizado por un objeto de $\cC$, y lo que nos dice
el diagrama a grandes rasgos es que el morfismo <<se comporta igual>> para
cualquier valor del parámetro.

\begin{example}\label{ex:nattrans}\;
  \begin{enumerate}
  \item En $\bVecf$, un isomorfismo de un espacio $V$ a su dual $V^*$ viene dado
    por una forma bilineal no degenerada, pero como esta no es única y no hay
    una forma general de elegir una, el isomorfismo no puede ser natural. Sin
    embargo, $\phi_V:V\to V^{**}$ dado por $\phi_V(x)(f)\coloneqq f(x)$ es un
    isomorfismo y se puede definir del mismo modo para todo $V$, por lo que
    $\phi$ es una transformación natural $\phi:1_{\bVecf}\to(^{**})$ y, de
    hecho, es una transformación natural $\phi:1_{K\dash\bVec}\to(^{**})$, donde
    el funtor $(^{**})$ viene dado por la composición
    $(^{**}):\bVec\overset{(^*)}{\to}\dual{\bVec}\overset{(^*)}{\to}\bVec$.
  \item Para cada $n\in\sNat$, el determinante de matrices $n\times n$ es una
    transformación natural. Más concretamente, si para un anillo conmutativo $R$
    tomamos la función determinante $\det_R:{\mathcal{M}}_n(R)\to R$, $\det_R$
    no es un homomorfismo de anillos porque no conserva la suma, pero sí es un
    homomorfismo de grupos multiplicativos $\det_R:\text{GL}_n(R)\to R^*$. Si
    $f$ es un homomorfismo de anillos, podemos definir $\text{GL}_n(f)$
    componente a componente y $f^*=f$, y entonces $\det:\text{GL}_n\to(\cdot)^*$
    es una transformación natural entre funtores $\bCRng\to\bGrp$.
  \end{enumerate}
\end{example}

\section{Categorías de funtores}

Las transformaciones naturales se pueden componer de varias formas. Una forma
sencilla es componiendo los morfismos imagen de dos transformaciones naturales
\[\cC\nattwos{\sigma}{\tau}\cD.\]

\begin{definition}
  Dados tres funtores $R,S,T:\cC\to\cD$ y dos transformaciones naturales
  $\sigma:R\to S$ y $\tau:S\to T$, llamamos \conc{composición vertical} de
  $\sigma$ y $\tau$ a la transformación natural $\tau\cdot\sigma:R\to T$ dada
  por $(\tau\cdot\sigma)_c\coloneqq \tau_c\circ \sigma_c$.
\end{definition}

Claramente esta composición es asociativa y tiene identidad, lo que sugiere la
siguiente definición.

\begin{definition}
  Dadas dos categorías $\cC$ y $\cD$, la \conc{categoría de funtores} de $\cC$ y
  $\cD$, $\cD^\cC$, es aquella que tiene como objetos los funtores $\cC\to\cD$,
  como morfismos las transformaciones naturales entre ellos, como composición la
  composición vertical y como identidad la \conc{transformación natural
    identidad}, que para un funtor $T:\cC\to\cD$ viene dada por
  $1_T(c)\coloneqq 1_{Tc}$.
\end{definition}

\begin{example}\;
  \begin{enumerate}
  \item Si $I$ es un conjunto, la categoría de funtores $\cC^I$ es precisamente
    la $I$-ésima potencia de $\cC$ en una categoría de categorías apropiada. En
    concreto, si $X$ es otro conjunto, $X^I$ es el conjunto de familias
    $(x_i)_{i\in I}$ con entradas en $X$, o de funciones $I\to X$.
  \item Si $X$ es un conjunto, $\{0,1\}^X$ es su conjunto potencia.
  \item $\cC^\bOne$ es isomorfa a $\cC$, mientras que $\cC^\bZero$ es la categoría unipuntual.
  \item $\cC^\bDown$ es la \conc{categoría de flechas} de $\cC$, cuyos objetos
    son los morfismos de $\cC$ y cuyos morfismos $f\to g$ son los pares de
    flechas $(h,k)$ para los que la figura \ref{fig:arrow-cat} conmuta.
    \begin{figure}
      \centering
      \begin{diagram}
        \path (0,2) node(A){$\cdot$} (2,2) node(AP){$\cdot$};
        \path (0,0) node(B){$\cdot$} (2,0) node(BP){$\cdot$};
        \draw[->] (A) -- node[left]{$f$} (B);
        \draw[->] (AP) -- node[right]{$g$} (BP);
        \draw[->] (A) -- node[above]{$h$} (AP);
        \draw[->] (B) -- node[below]{$k$} (BP);
      \end{diagram}
      \caption[Categoría de flechas.]{Morfismos $f\to g$ de la categoría de flechas.}
      \label{fig:arrow-cat}
    \end{figure}
  \item Si $M$ es un monoide, $\bSet^M$ es la categoría de las acciones de $M$
    sobre algún conjunto.
  \item Consideremos el \conc{funtor diagonal} $\Delta:\cC\to\cC^\cS$. Para cada
    objeto $c$, $\Delta c:\cS\to\cC$ es el funtor constante que lleva cada
    objeto a $c$ y cada morfismo a $1_c$, y para cada morfismo $f$,
    $\Delta f:\Delta a\to\Delta b$ es la transformación natural que lleva todos
    los objetos a $f$. Si $D:\cS\to\cC$ es un diagrama, un morfismo
    $\Delta c\to D$ es una fuente que conmuta con $D$, y un morfismo
    $D\to\Delta c$ es un sumidero que conmuta con $D$.
  \end{enumerate}
\end{example}

\begin{proposition}
  Dadas dos categorías $\cC$ y $\cD$:
  \begin{enumerate}
  \item Si $\cC$ y $\cD$ son pequeñas, también lo es $\cD^\cC$.
    \begin{proof}
      Los objetos de $\cD^\cC$ son funciones $\Mor{\cC}\to\Mor{\cD}$, los
      morfismos son funciones $\Ob{\cC}\to\Mor{\cD}$, y ambos conjuntos de
      funciones son pequeños.
    \end{proof}
  \item Si $\cC$ es pequeña y $\cD$ es una clase, $\cD^\cC$ es una clase.
    \begin{proof}
      Si $T:\cC\to\cD$ es un funtor, como $\Mor{\cC}$ es pequeño, su imagen por
      $\cC$ también y $T$ es pequeño, por lo que el conjunto de funtores es una
      clase, y análogamente para el de transformaciones naturales.
    \end{proof}
  \item Si $\cC$ es pequeña y $\cD$ tiene conjuntos hom pequeños, $\cD^\cC$
    tiene conjuntos hom pequeños.
    \begin{proof}
      Los elementos de $\hom_{\cD^\cC}(S,T)$ son funciones de $\Ob{\cC}$ a
      $\bigcup_{c\in\Ob{\cC}}\hom(Sc,Tc)$, pero estos conjuntos son pequeños y
      por tanto el conjunto de estas funciones también.
    \end{proof}
  \end{enumerate}
\end{proposition}

\section{Isomorfismos naturales}

Los isomorfismos de las categorías de funtores son particularmente importantes,
pues proporcionan una forma general de obtener isomorfismos que, además,
conmutan con los morfismos imágenes de los correspondientes funtores,
proporcionando una forma directa de pasar de uno al otro.

\begin{definition}
  Un \conc{isomorfismo natural} es un isomorfismo en una categoría de funtores,
  es decir, una transformación natural $\tau:S\to T$ en que todos los $\tau_c$
  son isomorfismos, y si existe decimos que los funtores $S$ y $T$ son
  \conc{naturalmente isomorfos}.
\end{definition}

\begin{example}\;
  \begin{enumerate}
  \item La transformación natural $\phi:1\to(^{**})$ del ejemplo \ref{ex:nattrans}
    no es un isomorfismo natural, pero sí lo es si restringimos el dominio y
    codominio de los funtores a $K\dash\bVecf$.
  \item Si $a$ y $b$ son objetos de una cierta categoría con coproducto
    $a\oplus b$, existe una biyección
    $\hom(a,d)\times\hom(b,d)\cong\hom(a\oplus b,d)$ natural respecto a $d$,
    como se muestra en la figura \ref{fig:nat-coproduct}.
    \begin{figure}
      \centering
      \begin{subfigure}{.45\textwidth}
        \begin{diagram}\smaller
          \path (0,2) node(ABX){$\hom(a,x)\times\hom(b,x)$} (4,2) node(SX){$\hom(a\oplus b,x)$};
          \path (0,0) node(ABY){$\hom(a,y)\times\hom(b,y)$} (4,0) node(SY){$\hom(a\oplus b,y)$};
          \draw[->] (ABX) -- node[right]{$\hom(a,f)\times\hom(b,f)$} (ABY);
          \draw[->] (SX) -- node[right]{$\hom(a\oplus b,f)$} (SY);
          \draw[->] (ABX) -- node[above]{$\tau_x$} (SX);
          \draw[->] (ABY) -- node[below]{$\tau_y$} (SY);
        \end{diagram}
        \caption{Diagrama asociativo.}
      \end{subfigure}
      \hfill
      \begin{subfigure}{.45\textwidth}
        \begin{diagram}\smaller
          \path (0,2) node(ABX){$(h,k)$} (4,2) node(SX){$h\oplus k$};
          \path (0,0) node(ABY){$(f\circ h,f\circ k)$}
                (4,0) node(SY){$(f\circ h)\oplus(f\circ k)=f\circ(h\oplus k)$};
          \tikzsquig (ABX) -- (ABY);
          \tikzsquig (SX) -- (SY);
          \tikzsquig (ABX) -- (SX);
          \tikzsquig (ABY) -- node[below]{\vphantom{$\tau_y$}} (SY);
        \end{diagram}
        \caption{Persecución de flechas.}
      \end{subfigure}
      \caption{Transformación natural en un coproducto.}
      \label{fig:nat-coproduct}
    \end{figure}
    Esto se puede generalizar de forma obvia a coproductos de una familia
    arbitraria de objetos.
  \item Del mismo modo, si existe el producto de los objetos $a$ y $b$, existe
    una biyección $\hom(c,a)\times\hom(c,b)\cong\hom(c,a\times b)$ natural
    respecto a $c$, que también se puede extender a productos arbitrarios.
  \end{enumerate}
\end{example}

Ciertas propiedades de funtores se pueden caracterizar mediante isomorfismos
naturales.

\begin{proposition}
  Un funtor $T:\cC\to\cD$ es una equivalencia si y sólo si existe $S:\cD\to\cC$
  tal que $S\circ T=1_\cC$ y $T\circ S=1_\cD$.
\end{proposition}
\begin{proof}
  Si $T$ es una equivalencia, hemos visto en la prueba de la proposición
  \ref{prop:cat-equiv} que existen $S:\cD\to\cC$ y un isomorfismo natural
  $h:T\circ S\to 1_\cD$. Además, si $c$ es un objeto de $\cC$,
  $h_{Tc}^{-1}:Tc\to TSTc$ es un isomorfismo, y como $T$ es fiel y pleno, existe
  un único $\mu_c:c\to STc$ con $T\mu_c=h_{Tc}^{-1}$. La naturalidad de
  $\mu:1_\cC\to S\circ T$ se deriva de la de $h$ y la fidelidad de $T$.

  Recíprocamente, si existen $S:\cD\to\cC$ e isomorfismos naturales
  $\mu:1_\cC\to S\circ T$ y $\sigma:T\circ S\to1_\cD$, para cada objeto $d$ en
  $\cD$, $T(Sd)\cong d$, y queda ver que $T$ es fiel y pleno. Es fiel porque,
  dados dos morfismos $f,g:a\to b$ en $\cC$,
  \[
    Tf=Tg\implies\mu_b\circ f=STf\circ\mu_a=STg\circ\mu_a=\mu_b\circ g\implies
    f=g,
  \]
  y es plena porque, dado un morfismo $g:Ta\to Tb$ en $\cD$,
  $f\coloneqq\mu_b^{-1}\circ Sg\circ\mu_a$ cumple $Tf=g$.
\end{proof}

\begin{definition}
  Un funtor $T:\cC\to\bSet$ es \conc{representable} por un objeto $c$ si es
  naturalmente isomorfo a $\hom(c,-)$.
\end{definition}

\begin{example}\;
  \begin{enumerate}
  \item Un funtor olvidadizo es representable por un objeto $c$ si y sólo si $c$
    es un objeto libre sobre $\{*\}$.
  \item El funtor olvidadizo de $\bBan$ no es representable.
  \end{enumerate}
\end{example}

\begin{proposition}
  Dados un funtor $T:\cC\to\bSet$, un objeto $c$ de $\cC$ y $x\in Tc$, existe
  una única transformación natural $\tau:\hom(c,-)\to T$ con $\tau_c(1_c)=x$.
\end{proposition}
\begin{proof}
  Si $\tau$ es tal transformación natural, $b$ es un objeto de $\cC$ y
  $f:c\to b$, necesariamente
  $\tau_b(f)=\tau_b(f\circ
  1_c)=(\tau_b\circ\hom(c,f))(1_c)=(Tf\circ\tau_c)(1_c)=(Tf)(x)$. Esta fórmula
  define una transformación natural, que es pues la única que cumple la
  condición.
\end{proof}

\section{Composición horizontal}

Además de componer transformaciones naturales verticalmente, podemos componer
dos transformaciones naturales $\cB\nats{\tau}\cC\nats{\sigma}\cD$ de la
siguiente manera.

\begin{definition}
  Dados cuatro funtores $S,T:\cB\to\cC$ y $S',T':\cC\to\cD$ y dos
  transformaciones naturales $\tau:S\to T$ y $\sigma:S'\to T'$, llamamos
  \conc{composición horizontal} de $\tau$ y $\sigma$ a la transformación natural
  $\sigma\circ\tau:S'\circ S\to T'\circ T$ dada por la figura
  \ref{fig:horiz-comp} como
  $(\sigma\circ\tau)_b\coloneqq T'\tau_b\circ\sigma_{Sb}=\sigma_{Tb}\circ
  S'\tau_b$, donde la conmutatividad se debe a la naturalidad de $\sigma$.
  \begin{figure}
    \centering
    \begin{diagram}
      \path (0,2) node(SS){$S'Sb$} (2,2) node(TS){$T'Sb$};
      \path (0,0) node(ST){$S'Tb$} (2,0) node(TT){$T'Tb$};
      \draw[->] (SS) -- node[left]{$S'\tau_b$} (ST);
      \draw[->] (TS) -- node[right]{$T'\tau_b$} (TT);
      \draw[->] (SS) -- node[above]{$\sigma_{Sb}$} (TS);
      \draw[->] (ST) -- node[below]{$\sigma_{Tb}$} (TT);
      \draw[->] (SS) -- node[above]{$\sigma\circ\tau$} (TT);
    \end{diagram}
    \caption{Composición horizontal de transformaciones naturales.}
    \label{fig:horiz-comp}
  \end{figure}
\end{definition}

Es fácil ver que la composición horizontal de transformaciones naturales es otra
transformación natural y que esta operación es asociativa. Además, la
transformación natural identidad en un funtor identidad es una identidad de esta
composición, por lo que las transformaciones naturales son los morfismos de una
categoría cuyos objetos son categorías.

La siguiente proposición es fácil de probar.

\begin{proposition}[Ley del intercambio]
  Dadas las transformaciones naturales
  \[
    \cB\nattwos{\sigma}{\tau}\cC\nattwos{\sigma'}{\tau'}\cD,
  \]
  se tiene $(\tau'\cdot\sigma')\circ(\tau\cdot\sigma)=(\tau'\circ\tau)\cdot(\sigma'\circ\sigma)$.
\end{proposition}
% \begin{proof}
%   La figura \ref{fig:nat-interchange} conmuta.
%   \begin{figure}
%     \centering
%     \begin{diagram}
%       \path (0,4) node(SS){$S'Sb$} (2,4) node(TS){$T'Sb$} (4,4) node(US){$U'Sb$};
%       \path (2,2) node(TT){$T'Tb$} (4,2) node(UT){$U'Tb$};
%       \path (4,0) node(UU){$U'Ub$};
%       \draw[->] (SS) -- node[above]{$\sigma'_{Sb}$} (TS);
%       \draw[->] (TS) -- node[above]{$\tau'_{Sb}$} (US);
%       \draw[->] (US) -- node[right]{$U'\sigma_b$} (UT);
%       \draw[->] (UT) -- node[right]{$U'\tau_b$} (UU);
%       \draw[->] (TS) -- node[right]{$T'\sigma_b$} (TT);
%       \draw[->] (TT) -- node[above]{$\tau'_{Tb}$} (UT);
%       \draw[->] (SS) -- node[above]{$\sigma'\circ\sigma$} (TT);
%       \draw[->] (TT) -- node[above]{$\tau'\circ\tau$} (UU);
%       \path (0,2) node(ST){$S'Tb$} (0,0) node(SU){$S'Ub$} (2,0) node(TU){$T'Ub$};
%       \draw[->] (SS) -- node[left]{$S'\sigma_b$} (ST);
%       \draw[->] (ST) -- node[left]{$S'\tau_b$} (SU);
%       \draw[->] (SU) -- node[below]{$\sigma'_{Ub}$} (TU);
%       \draw[->] (TU) -- node[below]{$\tau'_{Ub}$} (UU);
%       \draw[->] (ST) -- node[below]{$\sigma'_{Tb}$} (TT);
%       \draw[->] (TT) -- node[left]{$T'\tau_b$} (TU);
%     \end{diagram}
%     \caption{Ley del intercambio.}
%     \label{fig:nat-interchange}
%   \end{figure}
% \end{proof}

Esto muestra un ejemplo de \conc{bicategoría} o \conc{categoría bidimensional},
un par de categorías con el mismo conjunto de morfismos que cumple la ley del
intercambio y en que las identidades de una de las dos lo son de la otra. Las
categorías implicadas son, por supuesto, una categoría de categorías con sus
transformaciones naturales y la unión disjunta de las categorías de funtores
entre ellas.

%%% Local Variables:
%%% mode: latex
%%% TeX-master: "main"
%%% End:
